\chapter{Future Directions}

\section{Measuring the Impact of AI on Software Quality}

Measuring the impact of AI based tools on software quality remains an open challenge. Software quality is a complex concept that includes factors like maintainability, reliability, and security, which are often difficult to evaluate by themselves. Despite the existence of quality models, like ISO/IEC 25010, which provide a useful framework. Applying them consistently to AI assisted development workflows remains an ongoing research.

The evaluations of AI tools done now usually rely on productivity metrics or short-term studies. Even though these approaches can give an insight in specific capabilities, the could not be able to capture the long-term effects that AI assistance has on software quality. A great example we know is that, a greater development speed does not mean that the final product will be of higher quality, or even easier to maintain.

One more limitation is found in the lack of standardized evaluation methodologies. Different studies would focus on \textit{singular} aspects of quality, like code correctness or even bug detection. This makes comparing multiple tools' results a complicated task.

Future studies should include the development of more comprehensive evaluation frameworks. Researches involving development teams could provide insight on the influence AI-based tools have on software quality. These actions may help close the gap between controlled experiments and practical programming outcomes.

\section{Long-Term Maintainability and Technical Debt}

The effects that AI-tools have on software quality, long term, remain an open discussion. Even though these tools can improve a code structure, as well as the productivity of it, on a short-term basis, on a long-term on the influence they have over the structure and the constant evolution of programming is less understood. Maintainability, but also code readability, ease of modification and safety, could get affected by the assistance of AI.

A big concern is the accumulation of \textit{technical debt}. Generated code could prioritize functional over logical correctness, we can see this happen when developers rely on AI and adopt the suggestions given without properly understanding their significance. This will inevitably lead to codebases that are harder to debug, modernize, or refactor, increasing the need for maintenance over time.

On the other hand, these tools can also positively enhance the quality over long periods of time by encouraging standardized coding practices and by identifying quality-related issues early during the development phase. For example, automated detection of bugs or complex structures can help reducing the growth of technical debt.

Future studies should aim at figuring out the impact that AI-assisted programming has on extended period of time. Analyzing code evolution, cost of maintenance, and the developer comprehensions are all studies that could give insights on if AI-tools can mitigate the technical debt in real-word coding projects.

\section{Human Oversight and Trust in AI Tools}

With the increase of AI-based tools in development we see how important the concepts of human oversight and trust are. While AIs are able to provide assistance in tasks like code generation, analysis and testing, there is always a risk that their outputs are incorrect or even just suboptimal. As a result, relying on AI without evaluating the result given can lead to risking the software quality.

Trust in AI tools plays a significant role in the programming cycle. Excessively trusting such tools can often lead to the programming relying on them too much, often accepting generated suggestions without a proper check. On the other hand a lack of trust can reduce the benefits of having AI \textit{assist} the programming experience, limiting its effectiveness. A balance between trust and skepticism is incredibly important to achieve.

Human-in-the-loop methods are a promising directions that should be maintained. By having developers still actively involved in decision-making, it ensures that AI systems work as expected but most importantly as supportive tools rather than a replacement. Some features that AIs offer, and should be taken in consideration, like transparency over the thought process, explained suggestions, and customization levels of assistance may help programmers understand and evaluate better AI-generated recommendations.

Developments in the future should aim on making AI tools that cultivate meaningful human oversight and informed trust. Improving the explainability and providing a clearer feedback may help developers make more responsible decisions while also, enhancing the reliability of AI-assisted software systems.

\section{Ethical AI for Software Engineering}

The rise of AI in software engineering raises ethical questions beyond the technical performance. As these tools influence code generation, analysis, and testing, questions related to responsibility and transparency start to become more relevant. This means that Ethical AI in programming requires a careful consideration of how the tools are designed, deployed, and used practically.

One of the biggest concerns is the \textit{accountability}. When AI systems help software development, it could become unclear as to who's responsible for vulnerabilities, defects, or even unintended behaviors in the code. Despite the use of AI tools, responsibility remains with human developers, as well as with the companies, highlighting the importance for clear instructions and oversights.

Transparency and explainability also play an important role in ethical AI adoption. Programmers \textit{need} to understand, at least to some degree, why the AI tools return certain suggestions. Without a good enough transparency, deciding how reliable an AI-generated output or detecting biases from training data becomes complicated.

In the end, the ethical aspect is closely related to trust and long-term adoption. Tools that don't manage to respect privacy or intellectual property could stop being used, despite advantages that they may have over other tools. In the future, programmers that use AI should make their number one concern the ethical aspect of using such tools.
